\subsection{A property of vague convergence}

Lemma 5. --- Let X be a locally compact space, \(\mu\) a bounded positive measure on X, F a Banach space, and f a bounded function on X with values in F. The following conditions are equivalent:

\begin{enumerate}
\def\labelenumi{(\roman{enumi})}
\item
  The set of points of discontinuity of \(\mathbf{f}\) is \(\mu\) -negligible.
\item
  For every \(\varepsilon > 0\) , there exist elements \(\mathbf{a}_1, \ldots, \mathbf{a}_n\) of \(F\) , functions \(g_1, \ldots, g_n\) belonging to \(\mathcal{K}(X)\) , and a bounded continuous function \(h \geqslant 0\) on \(X\) such that \(|\mathbf{f} - g_1 \mathbf{a}_1 - \cdots - g_n \mathbf{a}_n| \leqslant h \leqslant 2 \sup |\mathbf{f}|\) on \(X\) , and \(\int h d\mu \leqslant \varepsilon\) . Denote by \(N\) the set of points of discontinuity of \(\mathbf{f}\) , and let \(M = \sup |\mathbf{f}|\) .
\setcounter{enumi}{0}
\item
  \(\Rightarrow\) (ii). Suppose that condition (i) is satisfied. Let \(\varepsilon > 0\) . The function f is \(\mu\) -integrable (No.~2, Cor. 4 of Prop. 5, and No.~6, Th. 5), therefore there exist \(a_{1}, \ldots, a_{n}\) in F and \(g_{1}, \ldots, g_{n}\) in \(\mathcal{K}(X)\) such that, on setting \(k = |f - g_{1}a_{1} - \cdots - g_{n}a_{n}|\) , we have \(\int k d\mu \leqslant \varepsilon/2\) (§3, No.~5,
Prop. 10). Multiplying \(g_1, \ldots, g_n\) by a suitable same element of \(\mathcal{K}(X)\) , we can suppose in addition that
\end{enumerate}

\[
\left| g _ {1} \mathbf {a} _ {1} + \dots + g _ {n} \mathbf {a} _ {n} \right| \leqslant \mathrm{M} = \sup | \mathbf {f} |
\]

on X, whence \(k \leqslant 2M\) . The set \(N'\) of points of discontinuity of \(k\) is contained in \(N\) , hence is negligible. For every \(x \in X\) , set \(l(x) = \limsup_{y \to x} k(y)\) .

Then \(2\mathrm{M} \geqslant l \geqslant k\) on X, and \(l = k\) on \(\mathrm{X} - \mathrm{N}'\) , that is, almost everywhere for \(\mu\) , therefore \(\int l d\mu \leqslant \varepsilon / 2\) . On the other hand, \(l\) is bounded and upper semi-continuous, hence is the lower envelope of the set of bounded continuous functions \(\geqslant l\) . Therefore there exists a bounded continuous function \(h \geqslant l\) on X such that \(h \leqslant 2\mathrm{M}\) and \(\int h d\mu \leqslant \int l d\mu + \varepsilon / 2\) (§4, No.~4, Cor. 2 of Prop. 5). Then \(\int h d\mu \leqslant \varepsilon\) and \(|\mathbf{f} - g_1\mathbf{a}_1 - \dots - g_n\mathbf{a}_n| \leqslant h\) .

\begin{enumerate}
\def\labelenumi{(\roman{enumi})}
\setcounter{enumi}{1}
\tightlist
\item
  \(\Rightarrow\) (i). Suppose that condition (ii) is satisfied. For every \(x \in X\) let \(\omega(x)\) be the oscillation of f at x (GT, IX, §2, No.~3). Let \(\varepsilon > 0\) . There exist \(a_{1}, \ldots, a_{n}, g_{1}, \ldots, g_{n}, h\) with the properties of (ii). For every \(x \in X\) , \(\omega(x)\) is the oscillation of \(f - g_{1}a_{1} - \cdots - g_{n}a_{n}\) at x, therefore \(\omega(x) \leqslant 2h(x)\) . Thus \(\int \omega d\mu \leqslant 2\varepsilon\) . Consequently, the set \(A_{\varepsilon}\) of \(x \in X\) such that \(\omega(x) \geqslant \sqrt{\varepsilon}\) satisfies \(\mu(A_{\varepsilon}) \leqslant 2\sqrt{\varepsilon}\) . This proves that \(\mu(N) \leqslant 2\sqrt{\varepsilon}\) , whence \(\mu(N) = 0\) .
\end{enumerate}

PROPOSITION 22. --- Let F be a Banach space, X a locally compact space, E the set of bounded positive measures on X, μ an element of E, and B a filter base on E. Assume that B converges vaguely to μ and that \textbar\textbar ν\textbar\textbar{} tends to \textbar\textbar μ\textbar\textbar{} with respect to B. Let f be a mapping of X into F satisfying the following conditions:

\begin{enumerate}
\def\labelenumi{(\roman{enumi})}
\item
  \(\mathbf{f}\) is bounded, and is integrable for \(\mu\) and for every measure that belongs to some element of \(\mathfrak{B}\) ;
\item
  the set of points of discontinuity of \(\mathbf{f}\) is \(\mu\) -negligible.
\end{enumerate}

Then \(\int \mathbf{f}d\nu\) tends to \(\int \mathbf{f}d\mu\) with respect to \(\mathfrak{B}\) .

Let \(\varepsilon > 0\) . There exist elements \(\mathbf{a}_1, \ldots, \mathbf{a}_n\) in \(F\) , functions \(g_1, \ldots, g_n\) in \(\mathcal{K}(X)\) , and a bounded continuous function \(h \geqslant 0\) on \(X\) , such that

\[
\left| \mathbf {f} - g _ {1} \mathbf {a} _ {1} - \dots - g _ {n} \mathbf {a} _ {n} \right| \leqslant h \leqslant 2 \sup | \mathbf {f} |
\]

on X and \(\int h d\mu \leqslant \varepsilon\) (Lemma 5). Let \(M = \sup |f|\) . There exist a compact subset K of X such that \(\mu^{*}(X - K) \leqslant \varepsilon\) (§4, No.~7, Prop. 12 and No.~6, Th. 4), a compact neighborhood \(K'\) of K in X, and a continuous mapping \(h'\) of X into {[}0,2M{]} such that \(h' = h\) on K, \(h' = 2M\) on \(X - K'\) ; replacing \(h'\) by \(\sup(h, h')\) , we can suppose in addition that \(h' \geqslant h\) . Then \(\int(h' - h) d\mu \leqslant 2M\mu^{*}(X - K) \leqslant 2M\varepsilon\) . On the other hand, \(h' = h_1 + 2M\) , where \(h_1 \in \mathcal{K}(X)\) . Taking into account §4, No.~7, Prop. 12, the number \(\int h^{\prime}d\nu = \int h_{1}d\nu +2\mathrm{M}\| \nu \| \text{ tends to } \int h_{1}d\mu +2\mathrm{M}\| \mu \| = \int h^{\prime}d\mu \text{ with respect to } \mathfrak{B}.\) There then exists an \(A\in \mathfrak{B}\) such that, for every \(\nu \in A\) ,

\[
\begin{array}{c} \Big | \int (g _ {1} \mathbf {a} _ {1} + \dots + g _ {n} \mathbf {a} _ {n}) d \nu - \int (g _ {1} \mathbf {a} _ {1} + \dots + g _ {n} \mathbf {a} _ {n}) d \mu \Big | \leqslant \varepsilon , \\ \int h d \nu \leqslant \int h ^ {\prime} d \nu \leqslant \int h ^ {\prime} d \mu + \varepsilon \leqslant \int h d \mu + 2 \mathrm{M} \varepsilon + \varepsilon \leqslant 2 (\mathrm{M} + 1) \varepsilon . \end{array}
\]

These inequalities imply

\[
\begin{array}{c} \Big | \int \mathbf {f}   d \nu - \int \mathbf {f}   d \mu \Big | \leqslant \\ \int h   d \nu + \Big | \int (g _ {1} \mathbf {a} _ {1} + \dots + g _ {n} \mathbf {a} _ {n})   d \nu - \int (g _ {1} \mathbf {a} _ {1} + \dots + g _ {n} \mathbf {a} _ {n})   d \mu \Big | + \int h   d \mu \\ \leqslant 2 (\mathrm{M} + 2) \varepsilon , \end{array}
\]

which proves the proposition.

Remark. --- The conditions (i) and (ii) of Proposition 22 are satisfied if f is continuous and bounded.

Example. --- Let us take for X the compact space U of complex numbers of absolute value 1. On setting \(\mu(f) = \int_{0}^{1} f(e^{2i\pi t}) dt\) for every \(f \in \mathcal{K}(U)\) , one defines a positive measure of mass 1 on U. On the other hand, let \(\theta\) be a real number; for every integer \(n \geqslant 0\) , let \(\nu_{n}\) be the unit mass placed at the point \(e^{2i\pi n\theta}\) of U, and let

\[
\mu_ {n} = \frac {1}{n + 1} (\nu_ {0} + \dots + \nu_ {n}),
\]

so that \(\mu_n\) is a positive measure of mass 1 on \(\mathbf{U}\) . Then, if \(\theta\) is irrational, \(\mu_n\) tends vaguely to \(\mu\) . For, since the linear combinations of the functions \(z \mapsto z^k\) ( \(k \in \mathbf{Z}\) ) are dense in \(\mathcal{K}(\mathbf{U})\) (GT, X, §4, No.~4, Prop. 8), it suffices to prove that \(\mu_n(z^k)\) tends to \(\mu(z^k)\) for \(k \in \mathbf{Z}\) . Now, for \(k = 0\) , \(\mu_n(z^k) = \mu(z^k) = 1\) ; for \(k \neq 0\) ,

\[
\mu_ {n} (z ^ {k}) = \frac {1}{n + 1} (1 + e ^ {2 i \pi k \theta} + e ^ {4 i \pi k \theta} + \dots + e ^ {2 i \pi k n \theta}).
\]

Since \(e^{2i\pi k\theta} \neq 1\) (because \(\theta\) is irrational), we infer that

\[
\left| \mu_ {n} (z ^ {k}) \right| = \left| \frac {1}{n + 1} \frac {e ^ {2 i \pi k (n + 1) \theta} - 1}{e ^ {2 i \pi k \theta} - 1} \right| \leqslant \frac {1}{n + 1} \frac {2}{\left| e ^ {2 i \pi k \theta} - 1 \right|},
\]

therefore \(\mu_n(z^k)\) tends to \(0 = \mu (z^k)\) . Under these conditions, Proposition 22 can be applied, and we see in particular that if A is a subset of U with negligible boundary with respect to \(\mu\) , then \(\mu_n(A)\) tends to \(\mu(A)\) . In other words, if \(p_n\) denotes the number of integers \(k \in [0, n]\) such that \(e^{2i\pi k\theta} \in A\) , then \(n^{-1}p_n\) tends to \(\mu(A)\) as \(n\) tends to \(+\infty\) .

